\section{Tecnologie utilizzate}
\label{sec:tecnologie}

% Per ogni tecnologia, spiega brevemente COSA è e PERCHÉ l'hai scelta
% (non limitarti a elencarla: motiva la scelta rispetto alle alternative).

\subsection{Linguaggi di programmazione}

\begin{itemize}[label={--}, leftmargin=1.2cm]
    \item \textbf{Kotlin}: utilizzato per lo sviluppo dell'applicazione Android.

    \item \textbf{Python}: linguaggio utilizzato per il backend tramite il framework Flask.
\end{itemize}

\subsection{Framework e librerie}

\begin{itemize}[label={--}, leftmargin=1.2cm]
    \item \textbf{Jetpack Compose}: utilizzato per costruire l'interfaccia.

    \item \textbf{Retrofit}: utilizzato per effettuare le chiamate REST al backend Flask.

    \item \textbf{kotlinx.serialization}: utilizzato per il parsing delle risposte JSON del backend.

    \item \textbf{Room}: per la persistenza locale dei dati.

    \item \textbf{WorkManager}: per la sincronizzazione periodica dei dati locali non ancora inviati al backend.

    \item \textbf{Wi-Fi Direct API}: usato per instaurare il collegamento P2P tra i dispositivi.

    \item \textbf{osmdroid}: usata per la visualizzazione della posizione in cui è stata effettuata la transazione all'interno della history.

    \item \textbf{Flask}: framework python usato per la creazione del webservice REST.

    \item \textbf{flask-jwt-extended}: estensione Flask per la gestione dell'autenticazione tramite JSON Web Token (JWT).

    \item \textbf{pymongo}: utilizzato dal backend per leggere e scrivere i dati nel database MongoDB.
\end{itemize}

\subsection{Strumenti di sviluppo}

\begin{itemize}[label={--}, leftmargin=1.2cm]
    \item \textbf{Android Studio}: IDE per lo sviluppo Android.

    \item \textbf{Docker Desktop}: utilizzato per avviare i container del backend.

    \item \textbf{Git / GitHub}: utilizzato per il versionamento del codice.

    \item \textbf{MongoDB Compass}: client di MongoDB per visualizzare il contenuto del database.
\end{itemize}
